\documentclass[11pt]{article}
\usepackage[margin=1.1in]{geometry}
\usepackage{amsmath,amssymb,amsthm,booktabs,microtype}
\usepackage[colorlinks=true,linkcolor=black,citecolor=black,urlcolor=blue]{hyperref}
\newtheorem{proposition}{Proposition}
\newtheorem{corollary}{Corollary}
\theoremstyle{remark}\newtheorem{remark}{Remark}
\title{\vspace{-2em}\textbf{Affine Term-Structure Decompositions Cannot Identify\\ the Source of Announcement-Window Yield Movements}}
\author{Research note}
\date{11 August 2026}

\begin{document}
\maketitle
\vspace{-1.5em}

\begin{abstract}
\noindent A growing literature decomposes the movement of long-term Treasury yields in windows around FOMC announcements into an expectations component and a term premium, and draws economic conclusions from the split. This note shows that in the Adrian--Crump--Moench (ACM) framework, and in any affine model with a fixed state vector spanned by the yield curve, the decomposition is a \emph{deterministic linear function of the contemporaneous curve movement}. It therefore contains no information beyond the shape of the yield response, and cannot serve as evidence about the economic source of that response. We further show that such a model misattributes a fixed fraction of a genuine permanent shift in the equilibrium rate to the term premium---empirically about 47\%---because its expectations component is, by construction, asymptotically constant. Both results are verified on the published ACM series.
\end{abstract}

\section{Setup}

Let $X_t\in\mathbb{R}^K$ be the pricing state, the first $K$ principal components of yields, following a stationary Gaussian VAR(1)
\begin{equation}
X_{t+1}=\mu+\Phi X_t+v_{t+1},\qquad v_{t+1}\sim N(0,\Sigma),\qquad \rho(\Phi)<1,
\end{equation}
with the short rate affine in the state, $r_t=\delta_0+\delta_1'X_t$. Under the essentially-affine specification of prices of risk, zero-coupon yields are affine,
\begin{equation}
y_t^{(n)}=A_n+B_n'X_t .
\end{equation}
Write $\bar X=(I-\Phi)^{-1}\mu$ for the unconditional mean, so that $\mathbb{E}_t[X_{t+h}]=\bar X+\Phi^h(X_t-\bar X)$.

The \emph{expectations} (risk-neutral) component of the $n$-period yield is the average expected future short rate under the physical measure,
\begin{equation}
rn_t^{(n)}\;=\;\frac1n\sum_{h=0}^{n-1}\mathbb{E}_t\!\left[r_{t+h}\right]
\;=\;a_n+b_n'X_t,
\end{equation}
where
\begin{equation}
b_n'\;=\;\delta_1'\,\Psi_n,\qquad
\Psi_n\;\equiv\;\frac1n(I-\Phi)^{-1}\!\left(I-\Phi^{\,n}\right),\qquad
a_n=\delta_0+\delta_1'\bar X-b_n'\bar X .
\label{eq:bn}
\end{equation}
The term premium is defined residually,
\begin{equation}
tp_t^{(n)}\;=\;y_t^{(n)}-rn_t^{(n)}\;=\;(A_n-a_n)+(B_n-b_n)'X_t .
\end{equation}

\section{The decomposition is a function of the curve}

\begin{proposition}\label{prop:span}
Let $\mathcal{N}=\{n_1,\dots,n_K\}$ be $K$ maturities such that $\mathbf{B}\equiv[B_{n_1},\dots,B_{n_K}]'$ is invertible, and let $y_t^{\mathcal N}$ collect the corresponding yields. Then for every maturity $n$,
\begin{equation}
\Delta rn_t^{(n)}=c_n'\,\Delta y_t^{\mathcal N},
\qquad
\Delta tp_t^{(n)}=\left(\iota_n-c_n\right)'\Delta y_t^{\mathcal N},
\qquad
c_n'\equiv b_n'\mathbf{B}^{-1},
\label{eq:cn}
\end{equation}
where $\iota_n$ selects $y^{(n)}$ from $y^{\mathcal N}$. The weights $c_n$ depend only on estimated parameters.
\end{proposition}

\begin{proof}
From (2), $y_t^{\mathcal N}=A^{\mathcal N}+\mathbf{B}X_t$, so $X_t=\mathbf{B}^{-1}(y_t^{\mathcal N}-A^{\mathcal N})$ and $\Delta X_t=\mathbf{B}^{-1}\Delta y_t^{\mathcal N}$. Substituting into $\Delta rn_t^{(n)}=b_n'\Delta X_t$ gives the first equality; the second follows from $\Delta tp=\Delta y-\Delta rn$.
\end{proof}

\begin{remark}
Proposition~\ref{prop:span} implies that a regression of $\Delta rn_t^{(n)}$ on $K$ contemporaneous yield changes has $R^2=1$ exactly. The decomposition is a fixed relabelling of curve movements, not an additional measurement.
\end{remark}

\begin{corollary}[Event studies inherit the map]\label{cor:event}
Consider an event study of a shock $z_t$,
\[
\Delta y_t^{(n)}=\alpha_n+\beta_n z_t+\varepsilon_{n,t},\qquad n\in\mathcal N\cup\{n\}.
\]
Then the estimated component responses satisfy, exactly,
\begin{equation}
\beta_n^{\,rn}=c_n'\beta^{\mathcal N},
\qquad
\beta_n^{\,tp}=\beta_n-c_n'\beta^{\mathcal N}.
\end{equation}
The split of an event-study response is a deterministic function of the estimated \emph{yield} response profile $\beta^{\mathcal N}$. Two shocks producing identical curve responses receive identical decompositions, whatever their economic sources. The decomposition therefore cannot constitute evidence about the source.
\end{corollary}

\section{The constant-endpoint bias}

\begin{proposition}[Misattribution of a permanent endpoint shift]\label{prop:endpoint}
Suppose the equilibrium short rate shifts permanently, immediately and credibly, $\delta_0\mapsto\delta_0+\kappa$. Then truly $\Delta rn^{(n)}=\kappa$ and $\Delta tp^{(n)}=0$ for every $n$, and the curve shifts in parallel, $\Delta y^{\mathcal N}=\kappa\boldsymbol{1}$. A model estimated with $(\delta,\Phi,\mu,\mathbf{B})$ held fixed reports
\begin{equation}
\widehat{\Delta rn}^{(n)}=\kappa\,(c_n'\boldsymbol{1}),
\qquad
\widehat{\Delta tp}^{(n)}=\kappa\left(1-c_n'\boldsymbol{1}\right).
\end{equation}
The attribution is correct if and only if $c_n'\boldsymbol{1}=1$.
\end{proposition}

\begin{proposition}[The expectations component is asymptotically deterministic]\label{prop:limit}
Since $\rho(\Phi)<1$, $\|\Phi^{\,n}\|\to0$ and hence from \eqref{eq:bn}
\begin{equation}
b_n=\frac1n\,(I-\Phi')^{-1}\!\left(I-\Phi^{\,n\prime}\right)\delta_1\;\longrightarrow\;0
\quad\text{as }n\to\infty,
\qquad
rn_t^{(n)}\;\longrightarrow\;\delta_0+\delta_1'\bar X .
\end{equation}
\end{proposition}

\noindent Proposition~\ref{prop:limit} is the structural statement behind Proposition~\ref{prop:endpoint}. The model's expectations component at long maturities converges to a constant, because the model assumes the long-run mean of the short rate is one. It has no channel through which a time-varying equilibrium rate can enter $rn$. Consequently \emph{any} true variation in the endpoint is orthogonal to the model's expectations component and is absorbed, in full, by the residual that the literature labels the term premium. A test of ``did the endpoint move?'' conducted through such a decomposition has no power against the alternative, and mechanically returns the answer ``term premium.''

\section{Empirical verification}

We verify both propositions on the published ACM daily series, June 1989 to August 2026.

\paragraph{Proposition~\ref{prop:span}.} Regressing $\Delta rn^{(10)}$ on daily changes in $\{1,2,3,5,7,10\}$-year yields gives $R^2=1.00000$. The fit is exact, as the proposition requires, and is unchanged across alternative spanning sets.

\paragraph{Corollary~\ref{cor:event}.} We estimate the response of the yield curve to a one-percentage-point revision in the FOMC's longer-run federal funds projection, in a window $[t-1,t+1]$ around each of 56 SEP meetings:
\[
\beta^{\mathcal N}=(+0.4,\;+2.9,\;+10.8,\;+22.4,\;+28.1,\;+31.8)\ \text{bp at }1,2,3,5,7,10\text{ years}.
\]
Applying \eqref{eq:cn} to this profile predicts the component responses. Table~\ref{tab:one} compares the prediction with the directly estimated component responses.

\begin{table}[h]\centering
\caption{Predicted versus estimated component responses (basis points per percentage point)}\label{tab:one}
\begin{tabular}{lcc}
\toprule
& Risk-neutral & Term premium\\
\midrule
Implied by $c_n'\beta^{\mathcal N}$ &$+2.5$ & $+29.3$\\
Directly estimated & $+2.6$ & $+29.2$\\
\bottomrule
\end{tabular}
\end{table}

\noindent The agreement is exact to estimation error. The finding that ``the Fed's longer-run projection moves the term premium and not expectations'' is, arithmetically, the statement that the front end of the curve did not move while the long end did.

\paragraph{Proposition~\ref{prop:endpoint}.} The empirical weights give $c_{10}'\boldsymbol{1}=0.529$, stable across spanning sets ($0.52$--$0.57$). A permanent, credible one-percentage-point rise in the equilibrium short rate---an event that is $100\%$ expectations by construction---is therefore reported by ACM as $+53$ bp of expectations and $+47$ bp of term premium.

\section{Implications}

\begin{enumerate}
\item Published decompositions of FOMC-window yield movements into expectations and term premia are re-expressions of the estimated curve response. They may be reported as descriptive summaries, but they cannot be cited as evidence for one economic mechanism over another.
\item Studies of the secular decline in long rates are the worst case for this bias, because the object of interest---a shifting equilibrium rate---is precisely the object the model assumes away. Proposition~\ref{prop:limit} implies the bias runs systematically toward the term premium.
\item Augmenting the state with an ``observed shifting endpoint'' $\tau_t$ resolves the bias \emph{only if} $\tau_t$ is measured independently of the curve. If $\tau_t=w'y_t$ for some loading $w$---for instance, a principal component of cumulative yield changes in announcement windows---then $rn_t$ remains affine in $y_t$, Proposition~\ref{prop:span} applies with modified weights, and the circularity is reinstated in a less visible form.
\end{enumerate}

\section{What restores identification}

Identification requires an object outside the span of the curve, or a restriction on dynamics rather than on the contemporaneous cross-section.

\paragraph{(a) Survey expectations.} Define $rn^{S,(n)}_t=\frac1n\sum_{h}\mathbb{E}^S_t[r_{t+h}]$ from surveys and $tp^{S,(n)}_t=y^{(n)}_t-rn^{S,(n)}_t$. Since $\mathbb{E}^S$ is measured independently of prices, Proposition~\ref{prop:span} does not apply. The cost is frequency: surveys are quarterly or semiannual, so daily event studies are unavailable and the method identifies levels rather than announcement responses.

\paragraph{(b) Dynamic restrictions.} Estimate the cumulative response by local projection,
\begin{equation}
y^{(n)}_{t+h}-y^{(n)}_{t-2}=\alpha_h+\theta_h z_t+u_{t+h},\qquad h=0,1,\dots,H .
\end{equation}
A permanent endpoint revision implies $\theta_h\to\kappa\neq0$; a transitory risk-premium movement implies $\theta_h\to0$. This uses only realised yields and imposes no affine mapping, so it is immune to Propositions~\ref{prop:span}--\ref{prop:limit}. It is, on the present argument, the only route that identifies the source of an announcement-window movement at daily frequency.

\paragraph{(c) Cross-sectional shape, read structurally.} The response profile $\beta^{\mathcal N}$ is itself model-free. An immediate, credible, permanent endpoint shift implies $\beta_n\approx\kappa$ for all $n$, including the shortest. The estimated profile above has $\beta_1=+0.4$ bp against $\beta_{10}=+31.8$ bp, which rejects that hypothesis without any decomposition. It does not, however, separate a term premium from a \emph{phased-in} expectations revision, since both load increasingly in maturity.

\vspace{1em}
\noindent\rule{\textwidth}{0.4pt}\\[0.4em]
\footnotesize
\textbf{Data.} Adrian, Crump \& Moench daily term-structure decomposition (\texttt{ACMY}, \texttt{ACMRNY}, \texttt{ACMTP}), June 1989--August 2026. FOMC longer-run projections from the Summary of Economic Projections dot-plot archive, 57 meetings, January 2012--March 2026; the right-hand-side variable is the revision in the cross-participant mean of the longer-run federal funds dot. Standard errors on the event study are Newey--West with four lags.

\vspace{0.6em}
\noindent\textbf{References.} Adrian, T., R. Crump and E. Moench (2013), ``Pricing the Term Structure with Linear Regressions,'' \emph{Journal of Financial Economics} 110(1), 110--138. \; Bauer, M. and G. Rudebusch (2020), ``Interest Rates under Falling Stars,'' \emph{American Economic Review} 110(5), 1316--1354. \; Crump, R., S. Eusepi and E. Moench (2018), ``The Term Structure of Expectations and Bond Yields,'' FRBNY Staff Report 775. \; Hillenbrand, S. (2025), ``The Fed and the Secular Decline in Interest Rates,'' \emph{Review of Financial Studies} 38(4), 981--1013. \; Kim, D. and J. Wright (2005), ``An Arbitrage-Free Three-Factor Term Structure Model,'' FEDS 2005-33.
\end{document}
